\usepackage{graphicx}%
\usepackage{multirow}%
\usepackage{amsmath,amssymb,amsfonts}%
\usepackage{amsthm}%
\usepackage{mathrsfs}%
\usepackage[title]{appendix}%
\usepackage{xcolor}%
\usepackage{textcomp}%
\usepackage{manyfoot}%
\usepackage{booktabs}%
\usepackage{multicol}
\usepackage{fancyvrb}
\usepackage{geometry}
\usepackage{titlesec}
\usepackage{xcolor}
\usepackage{parskip}
\usepackage{array}
\usepackage{multirow}
\usepackage{booktabs}
% Commented out biblatex because sn-jnl class loads natbib (incompatible)
% If you need bibliographies, use natbib commands: \citep{}, \citet{}, \bibliography{}
% \usepackage[backend=biber,style=apa]{biblatex}
% \addbibresource{references.bib}
\usepackage{hyperref}
\usepackage{xurl}             % Permite que los \url{} largos hagan salto de línea en cualquier caracter (evita overfull hbox en la sección de Referencias)
\usepackage{fontawesome5} % For the GitHub icon
\usepackage[justification=centering]{caption}
\usepackage[linesnumbered,ruled,vlined]{algorithm2e}
\setlength{\topmargin}{-0.4in}
\setlength{\textheight}{9in}
\usepackage{comment}
\usepackage{listings}%
% In your preamble:
\usepackage{enumitem} % For custom list formatting
\newlist{questions}{enumerate}{1}
\setlist[questions,1]{
  label=\textbf{Q\arabic*:}, % Label format "Q1:", "Q2:", ...
  leftmargin=2em,           % Indent spacing
  itemsep=1em               % Space between questions
}



%%%%



\raggedbottom



% Code Packages

\usepackage{listings}
\usepackage{xcolor}

\definecolor{headerbackground}{rgb}{1, 1, 1}
\definecolor{rulecolor}{rgb}{0.7, 0.7, 0.7}

% Define C keywords including standard library functions
\lstdefinestyle{algorithmstyle}{
  backgroundcolor=\color{white},
  basicstyle=\ttfamily\scriptsize,
  keywordstyle=\bfseries, % Bold for keywords
  commentstyle=\itshape, % Italic for comments
  stringstyle=,
  numberstyle=\tiny,
  numbers=left,
  numbersep=10pt,
  breaklines=true,
  frame=trbl,
  framerule=0.8pt,
  rulecolor=\color{rulecolor},
  framesep=5pt,
  xleftmargin=15pt,
  xrightmargin=5pt,
  language=C, % This is crucial
  tabsize=2,
  showstringspaces=false,
  alsoletter={\_},
  morekeywords={printf, scanf, malloc, free, sizeof, puts, gets}, % Add standard library functions
}

% ============================================ %
% Configuración adicional agregada para el      %
% reporte de CAN Bus (protocolo CAN)            %
% ============================================ %
\usepackage{siunitx}          % Unidades (Mbit/s, ns, V, etc.)
\sisetup{per-mode=symbol}
\usepackage{tabularx}         % Columnas que ajustan su ancho al \linewidth (evita tablas desbordadas)
\usepackage{ragged2e}         % \RaggedRight dentro de columnas X (mejor que justificado forzado)
\newcolumntype{Y}{>{\RaggedRight\arraybackslash}X}

% Resaltado rápido de bits dominantes/recesivos dentro del texto
\newcommand{\dominant}{\textbf{dominante (0)}}
\newcommand{\recessive}{\textbf{recesivo (1)}}

\newcommand{\codeheader}[1]{%
  \noindent%
  \fcolorbox{rulecolor}{headerbackground}{%
    \makebox[\dimexpr\linewidth-2\fboxsep-2\fboxrule][l]{%
      \textbf{Code:} \ttfamily{#1}%
    }%
  }%
  \vspace{-0.5\baselineskip}
}